\subsection{Podejścia oparte na promptowaniu (in-context learning)}
\label{subsec:llm_acr_prompting}

Najprostszy sposób na użycie LLM w automatycznym code review to potraktowanie modelu jako mechanizmu
"wnioskowania z kontekstu" (\emph{in-context learning}). W praktyce oznacza to, że wagi modelu nie są zmieniane, a
sterowanie odbywa się poprzez sam prompt: instrukcję, format odpowiedzi i kilka przykładów (albo ich brak - tzw. zero-shot).
Jakość wyników zależy głównie od tego, co i jak trafi do treści promptu - od doboru przykładów/demonstracji (\emph{few-shot}),
ale też od tego, ile kontekstu dołożymy i jak bardzo jest on czysty (np. czy nie ma szumu, sprzecznych wskazówek, niestandardowych 
formatów). Z punktu widzenia wdrożenia takich systemów jest to dobra wiadomość: nie jest potrzebna infrastruktura treningowa, 
możliwe jest szybkie iterowanie. Jednak literatura dość konsekwentnie pokazuje, że takie 
podejście bywa kapryśne i wrażliwe na detale, czasem wręcz na miniaturowe zmiany w promptcie.

Pornprasit i Tantithamthavorna opisują prompting jako alternatywę dla dostrajania, zwłaszcza w obliczu
problemów z danymi na starcie (tzw. \emph{cold-start}). Porównują kilka wariantów, m.in. zero-shot 
i few-shot, a także podejście z personą i bez persony. Ciekawy i praktyczny jest u nich wątek automatycznego doboru
przykładów: few-shot nie jest "ręcznie wybrane", tylko oparte o wyszukiwanie podobnych przypadków metodą BM25.
Jednocześnie liczba przykładów jest ograniczona do trzech, dzięki czemu określić to można jako świadomy kompromis
pomiędzy kosztem tokenów a stabilnością odpowiedzi \cite{Pornprasit2024FineTuningPromptingCR}. Wyniki autorów
zawierają również nieoczywiste wnioski: rozbudowane, wielokrokowe instrukcje nie zawsze pomagają, a w części ustawień lepiej 
działały prompty prostsze, ale o konsekwentnej strukturze \cite{Pornprasit2024FineTuningPromptingCR}. Dla projektu
narzędzia do automatycznego code review jest to dobry argument, aby budować szablony promptów
(np. dla każdego typu plik lub języka programowania) i testować je w kontrolowany sposób, zamiast dołączać co raz następne reguły
bo tak podpowiada logika.

Haider i in. patrzą na prompting w zadaniu generowania komentarzy recenzenckich (RCG) z szerszej perspektywy.
Autorzy zestawiają \emph{few-shot prompting} dla modeli zamkniętych z podejściami opartymi 
o dostrajanie (QLoRA) dla modeli otwartych. W samym wariancie promptingowym też pojawia 
się BM25 do wyboru przykładów, co wzmacnia wniosek, że w praktyce "promptowanie" 
często oznacza ustrukturyzowany ICL (\emph{in-context learning}) wsparty wyszukiwaniem, a nie tylko ręcznie napisany prompt 
\cite{Haider2024PromptingFineTuningRCG}. U Haider i in. ważny jest także kolejny szczegół: relacja między dodatkowymi metadanymi 
a jakością nie jest liniowa i jednoznaczna. Dokładanie np. grafu wywołań funkcji czy streszczenia 
kodu nie zawsze poprawiało rezultaty; efekt zależał między innymi od tego, jak 
dużo kontekstu model realnie potrafi sensownie wykorzystać \cite{Haider2024PromptingFineTuningRCG}.
Wniosek inżynierski na podstawie powyższych obserwacji to: wzbogacanie promptu trzeba traktować jako decyzję projektową i sprawdzać wpływ 
ablacją oraz pomiarem, zamiast przyjmować założenie "im więcej kontekstu, tym lepiej".

Kontekst w promptowaniu ma jeszcze jeden wymiar: nie tylko "kod", ale też 
opis zadania i intencja. Cihan i in. pokazują to na przykładzie procedury 
podobnej do realnego procesu recenzji zmian: model najpierw klasyfikuje poprawność kodu, a 
potem - dla przypadków uznanych za błędne - generuje poprawkę, którą weryfikuje 
się testami jednostkowymi. W ich ustawieniu dorzucenie opisu problemu do promptu poprawiało 
wyniki: dla zbioru 492 bloków kodu GPT-4o osiągnął 68.50\% trafności klasyfikacji poprawności 
oraz 67.83\% skuteczności poprawek, ale dalej występowało ryzyko regresji na poziomie 10.43\% 
\cite{Cihan2025EvaluatingLLMsForCodeReview}. To jest o tyle istotne, że nawet przy sensownie zaprojektowanym promptcie 
nie znika problem "złych zmian". W automatycznym \emph{code review} nie powinno się 
się tego zostawić bez zabezpieczeń: walidacja testami, analiza statyczna czy reguły blokujące 
wciąż wyglądają na konieczny element całego pipeline'u.

W tle tych prac przewija się też temat oceny jakości. Tufano i 
in. zwracają uwagę, że w zadaniach generatywnych (w tym w automatyzacji code 
review) łatwo wpaść w pułapkę metryk czysto ilościowych, a dodatkowo problemy w 
danych (np. błędne mapowania komentarz--zmiana) potrafią sztucznie zawyżać albo zaniżać wyniki \cite{Tufano2024CodeReviewAutomation}. 
Przy promptowaniu ma to szczególne znaczenie, bo iteracje są szybkie i często 
różnice są "na granicy": coś wygląda lepiej w jednej metryce, gorzej w 
innej, a jakościowo bywa różnie. W praktyce to sugeruje łączenie metryk automatycznych 
z oceną jakościową i testami możliwie bliskimi realnemu procesowi recenzji, a nie 
poleganie na jednym wskaźniku.

W efekcie prompting jest sensownym punktem startu dla automatycznego \emph{code review}, zwłaszcza 
gdy danych jest mało i trzeba szybko zbudować prototyp. Jednocześnie z literatury 
da się wyciągnąć kilka dość konkretnych zasad: przykłady warto dobierać automatycznie (np. 
BM25) i kontrolować ich liczbę; kontekst lepiej dobierać selektywnie i 
sprawdzać jego wpływ ablacją; a ocenę wyników oprzeć nie tylko o dopasowanie 
tekstu, ale też o mechanizmy, które łapią błędy merytoryczne i ryzyko regresji 
(testy, statyczna analiza, reguły). To nie rozwiązuje wszystkiego, ale przynajmniej ustawia prompting 
jako element większego procesu, a nie "magiczny prompt", który wykona recenzję sam 
z siebie.